\chapter{Generating Path Candidates}\labappendix{path_candidates}

As discussed in \refch{path-tracing} and \refappendix{tutorial}, deterministic (exact) ray tracing can be decomposed into two tasks: generating path candidates and then solving for the exact coordinates of each candidate. As a reminder, a path candidate is an ordered interaction sequence between a \glsxtrshort{tx} and an \glsxtrshort{rx}; geometric validity is verified only afterward (see \refsec{combinatorial-bottleneck}).

To optimize this computationally intensive process, modern differentiable ray tracers such as DiffeRT and Sionna~RT usually partition it into two stages:
\begin{enumerate}
  \item \emph{Candidate generation}, to identify ordered sequences of primitives (e.g., triangles) that a ray could potentially traverse.
  \item \emph{Path tracing}, to compute the exact coordinate path for each sequence.
\end{enumerate}

This appendix focuses only on the first stage: path candidate generation. In this formulation, listing candidates is equivalent to enumerating paths in a directed graph that connects scene objects with the \glsxtrshort{tx} and \glsxtrshort{rx} nodes, as illustrated in \reffig{path-candidates-graph} in \refch{path-tracing}.

\section{Scene with Known Visibility}

Let us consider again the simple 2D scene in \reffig{path-candidate-visibility-bis}. In preprocessed environments or simplified geometries, visibility between objects may be known or trivially computed.

\begin{figure}[htpb]
  \centering
  \inputtikz{path-candidate-visibility}
  \caption{2D scenario with triangular-shaped objects on which reflection or diffraction can occur, repeated here for convenience from \reffig{path-candidate-visibility}. Surfaces are colored black and edges are colored \textcolor{red}{red}.}
  \labfig{path-candidate-visibility-bis}
\end{figure}

This visibility is encoded as an adjacency matrix (visibility matrix), consistent with what we saw in \refch{efficient-path-tracing}. A value of one indicates an unobstructed \gls{line-of-sight} between two graph nodes.

\begin{figure}[htpb]
  \centering
  \begin{equation*}
    %\resizebox{.5\textwidth}{!}{
    \renewcommand{\arraystretch}{.6}
    %\setlength{\bigstrutjot}{2ex}
    \makeatletter\setlength\BA@colsep{3pt}\makeatother
    \begin{blockarray}{ccccccccccccccc}
      & \text{\glsxtrshort{tx}} & s_1 & s_2 & s_3 & s_4 & s_5 & s_6 & e_1 & e_2 & e_3 & e_4 & e_5 & e_6 & \text{\glsxtrshort{rx}}\\
      \begin{block}{c(>{\medspace}c|cccccc|cccccc|c<{\medspace})}
        \text{\glsxtrshort{tx}} & {} & 1 & {} & {} & {} & {} & 1 & 1 & 1 & {} & 1 & {} & 1 & 1\\
        \BAhhline{~~------------~}
        s_1 & {} & {} & {} & {} & {} & {} & \bsit{1} & {} & {} & {} & \bsit{1} & {} & & {} & {}\\ %s1
        s_2 & {} & {} & {} & {} & {} & {} & \bsit{1} & {} & {} & {} & \bsit{1} & {} & \bsit{1} & 1\\ %s2
        s_3 & {} & {} & {} & {} & {} & {} & {} & {} & {} & {} & {} & {} & {} & {}\\ %s3
        s_4 & {} & {} & {} & {} & {} & {} & {} & {} & {} & {} & {} & {} & {} & {}\\ %s4
        s_5 & {} & {} & {} & {} & {} & {} & {} & {} & {} & {} & {} & {} & {} & {}\\ %s5
        s_6 & {} & \bsit{1} & \bsit{1} & {} & {} & {} & {} & \bsit{1} & \bsit{1} & \bsit{1} & {} & {} & {} & 1\\ % s6
        e_1 & {} & {} & {} & {} & {} & {} & \bsit{\color{red}1} & {} & {} & {} & \bsit{\color{red}1} & {} & {} & {}\\ % e1
        e_2 & {} & {} & {} & {} & {} & {} & \bsit{\color{red}1} & {} & {} & {} & \bsit{\color{red}1} & {} & \bsit{\color{red}1} & 1\\ %e2
        e_3 & {} & {} & {} & {} & {} & {} & \bsit{\color{red}1} & {} & {} & {} & \bsit{\color{red}1} & {} & \bsit{\color{red}1} & 1\\ %e3
        e_4 & {} & \bsit{\color{red}1} & \bsit{\color{red}1} & {} & {} & {} & {} & \bsit{\color{red}1} & \bsit{\color{red}1} & \bsit{\color{red}1} & {} & {} & {} & 1 \\ %e4
        e_5 & {} & {} & {} & {} & {} & {} & {} & {} & {} & {} & {} & {} & {} & {}\\ %e5
        e_6 & {} & {} & \bsit{\color{red}1} & {} & {} & {} & {} & {} & \bsit{\color{red}1} & \bsit{\color{red}1} & {} & {} & {} & 1\\
        \BAhhline{~~------------~}
        \text{\glsxtrshort{rx}} & {} & {} & {} & {} & {} & {} & {} & {} & {} & {} & {} & {} & {} & {}\\
      \end{block}
    \end{blockarray}
    %}
  \end{equation*}
  \caption{Adjacency matrix, $\mathcal{G}$, generated from the scenario illustrated in \reffig{path-candidate-visibility-bis}, repeated here for convenience from \reffig{visibility-adjacency-matrix}. Each row of this $14 \times 14$ matrix refers to the visible objects as seen from the corresponding object. For readability, zeros are omitted. The black coefficients indicate a possible reflection, while the \textcolor{red}{red} ones indicate diffraction.}
  \labfig{visibility-adjacency-matrix-bis}
\end{figure}

If we omit the endpoint nodes (\glsxtrshort{tx} and \glsxtrshort{rx}), the matrix becomes symmetric because pairwise \gls{line-of-sight} visibility is bidirectional. Endpoint-specific visibility maps are shown in \reffig{path-candidate-visibility-from-bis}.

\begin{figure*}[htpb]
  \centering
  \begin{subfigure}{0.49\contentwidth}
    \centering
    \inputtikz{path-candidate-visibility-from-tx}
    \caption{Objects visible from \glsxtrshort{tx}.}
    \labfig{path-candidate-visibility-from-tx-bis}
  \end{subfigure}
  \begin{subfigure}{0.49\contentwidth}
    \centering
    \inputtikz{path-candidate-visibility-from-rx}
    \caption{Objects visible from \glsxtrshort{rx}.}
    \labfig{path-candidate-visibility-from-rx-bis}
  \end{subfigure}
  \caption{Example 2D visibility for the \glsxtrfull{tx} and \glsxtrfull{rx}, repeated here for convenience from \reffig{path-candidate-visibility-from}. Objects are faded if they are not visible from the corresponding node.}
  \labfig{path-candidate-visibility-from-bis}
\end{figure*}

Using a directed graph (\texttt{DiGraph} in DiffeRT's core library), we build the interaction graph from this adjacency matrix and enumerate candidates between the \glsxtrshort{tx} and the \glsxtrshort{rx}. Because the matrix is sparse, the candidate set is typically much smaller than in a complete-graph model.

\begin{listing*}[htpb]
  \caption{Constructing a directed graph from a known visibility matrix.}
  \lablisting{adjacency-matrix}
  \inputmintedmarker[]{python}{codes/wide/path-candidates.py}{adjacency_matrix}
\end{listing*}

\begin{listing}[htpb]
  \caption{Enumerating known-visibility path candidates with \texttt{DiGraph}.}
  \lablisting{graph-known-visibility}
  \inputmintedmarker[]{python}{codes/wide/path-candidates.py}{enumerate_paths_di_graph}
\end{listing}

The output from \reflisting{graph-known-visibility} is shown in \reflisting{graph-known-visibility-output}.
\begin{listing}[htpb]
  \caption{Output from \reflisting{graph-known-visibility}.}
  \lablisting{graph-known-visibility-output}
\begin{minted}{text}
#001: [ 0  1  6 13]
#002: [ 0  1 10 13]
#003: [ 0  6  2 13]
#004: [ 0  6  8 13]
#005: [ 0  6  9 13]
#006: [ 0  7  6 13]
#007: [ 0  7 10 13]
#008: [ 0  8  6 13]
#009: [ 0  8 10 13]
#010: [ 0  8 12 13]
#011: [ 0 10  2 13]
#012: [ 0 10  8 13]
#013: [ 0 10  9 13]
#014: [ 0 12  2 13]
#015: [ 0 12  8 13]
#016: [ 0 12  9 13]
\end{minted}
\end{listing}

\section{Scene with Unknown Visibility}

In many practical scenarios, exact visibility preprocessing is costly, especially in dense 3D environments (see \refch{efficient-path-tracing}).

Consequently, one may assume that every object can connect to every other object, i.e., a \emph{complete graph}. This is a worst-case candidate generator: it is simple and highly parallel, but it deliberately over-generates candidates that must later be rejected by geometric validity checks (\refsec{combinatorial-bottleneck}).

DiffeRT adopts this approach through the \texttt{CompleteGraph} class, which generates path indices an order of magnitude faster than an equivalent \texttt{DiGraph}. However, when modeling the \glsxtrshort{tx} and the \glsxtrshort{rx}, we generally do not want these terminal nodes to appear multiple times in a path: \glsxtrshort{tx} should only originate paths, and no node should connect back to \glsxtrshort{tx}. Similarly, \glsxtrshort{rx} should only terminate paths. To model this behavior with \texttt{CompleteGraph}, one must:
\begin{enumerate}
  \item Initialize a complete graph strictly \emph{without} the \glsxtrshort{tx} and the \glsxtrshort{rx} nodes.
  \item Generate path permutations programmatically (via \texttt{all\_paths}), specifying that the respective start and terminal indices (\texttt{from\_} and \texttt{to}) fall outside the graph's core logic.
\end{enumerate}
Because the implementation is optimized for complete graphs, it assumes that \texttt{from\_} is connected to all nodes in the graph, and symmetrically, that any node can connect to \texttt{to}.

\begin{listing}[htpb]
  \caption{Generating paths using \texttt{CompleteGraph} under unknown visibility assumptions.}
  \lablisting{graph-unknown-visibility}
  \inputmintedmarker[]{python}{codes/wide/path-candidates.py}{enumerate_paths_complete_graph}
\end{listing}

The (truncated) output from \reflisting{graph-unknown-visibility} is shown in \reflisting{graph-unknown-visibility-output}. Note that the candidate set is much larger than in the \texttt{DiGraph} case (132 candidates vs. 16).

\begin{listing}[htpb]
  \caption{Output from \reflisting{graph-unknown-visibility}.}
  \lablisting{graph-unknown-visibility-output}
\begin{minted}{text}
#001: [12  0  1 13]
#002: [12  0  2 13]
#003: [12  0  3 13]
#004: [12  0  4 13]
#005: [12  0  5 13]
#006: [12  0  6 13]
#007: [12  0  7 13]
#008: [12  0  8 13]
#009: [12  0  9 13]
#010: [12  0 10 13]
#011: [12  0 11 13]
#012: [12  1  0 13]
% ... 108 intermediate paths omitted ...
#121: [12 10 11 13]
#122: [12 11  0 13]
#123: [12 11  1 13]
#124: [12 11  2 13]
#125: [12 11  3 13]
#126: [12 11  4 13]
#127: [12 11  5 13]
#128: [12 11  6 13]
#129: [12 11  7 13]
#130: [12 11  8 13]
#131: [12 11  9 13]
#132: [12 11 10 13]
\end{minted}
\end{listing}

\section{Beware of the Number of Path Candidates}

The dominant limitation is combinatorial growth (\refsec{computational-complexity}, \refsec{exponential-path-generation}): the number of path candidates increases exponentially with the interaction order $K$.

For the complete-graph setting used here (with the \glsxtrshort{tx} and the \glsxtrshort{rx} handled as endpoints outside the core graph), the number of candidates for a fixed interaction order $K$ is
\begin{equation}
  N \times (N-1)^{K-1},
\end{equation}
where $N$ is the number of scene objects in the complete graph. This fixed-order term is consistent with the growth discussed in \refsec{computational-complexity}.

This induces the same trade-off emphasized in \refch{efficient-path-tracing}: a \texttt{DiGraph} with visibility pruning reduces the candidate count, while a \texttt{CompleteGraph} avoids preprocessing but incurs larger combinatorial growth.

Because full \texttt{CompleteGraph} implementations demand substantial memory allocations for higher interaction orders, tracing algorithms must either estimate a filtered visibility matrix---even if imperfect compared to mathematically exact ray tracing---or explicitly manage memory by iterating over candidate branches in smaller chunks. DiffeRT facilitates this through specialized block-processing methods (e.g., \texttt{*chunks\_iter} methods in \texttt{differt\_core.rt}), allowing continuous generation without catastrophic memory overflow.

Finally, to illustrate this scaling divergence, we compare the number of candidate paths produced by the directed graph (with exact visibility) and the complete graph across increasing depths.

\begin{listing*}[htpb]
  \caption{Comparing candidate counts between \texttt{DiGraph} and \texttt{CompleteGraph}.}
  \lablisting{differt-graph-comparison}
  \inputmintedmarker[]{python}{codes/wide/path-candidates.py}{count_paths}
\end{listing*}

The output of \reflisting{differt-graph-comparison} is shown in \reflisting{differt-graph-comparison-output}, demonstrating the exponential growth difference between the two graph models as depth increases.
\begin{listing*}[htpb]
  \caption{Output of \reflisting{differt-graph-comparison}.}
  \lablisting{differt-graph-comparison-output}
\begin{minted}{text}
depth = 2:      1 (DiGraph) vs      1 (CompleteGraph)
depth = 3:      4 (DiGraph) vs     12 (CompleteGraph)
depth = 4:     16 (DiGraph) vs    132 (CompleteGraph)
depth = 5:     56 (DiGraph) vs   1452 (CompleteGraph)
depth = 6:    192 (DiGraph) vs  15972 (CompleteGraph)
depth = 7:    680 (DiGraph) vs 175692 (CompleteGraph)
\end{minted}
\end{listing*}
